% GENERATED FILE. Edit canonical lessons, then run npm run generate:books.
% canonical-source-hash: fnv1a64:42ec7c1d800233f3
% canonical-lessons: KA-C486-hingaru, KA-C486-alikallu, KA-C486-mabbu, KA-C486-musuku, KA-C486-himapata

\chapter{Late Rains and Hail}
\label{ch:ka-late-rains-and-hail}
\hlchaptermodality{486}

\begin{chapteropening}
\textbf{By the end of this chapter:} \emph{I can say the late monsoon, hail, haze, to shroud and snowfall.}

Complete the last lesson of chapter 486: I can say the late monsoon, hail, haze, to shroud and snowfall.
\end{chapteropening}

\section[\texorpdfstring{\kn{ಹಿಂಗಾರು} (hiṅgāru)}{ಹಿಂಗಾರು (hiṅgāru)}]{\kn{ಹಿಂಗಾರು} (hiṅgāru) — the late monsoon; the retreating rains}

\label{lesson:KA-C486-hingaru}



\begin{quote}
Before the new one: say the Kannada for a downpour, then the Kannada for to get soaked.
\end{quote}

\subsection*{You'll want to know: \kn{ಹಿಂಗಾರು}}
\textbf{\kn{ಹಿಂಗಾರು}} — \emph{hiṅgāru} — \textquotedblleft{}the late monsoon; the retreating rains\textquotedblright{}.

\begin{quote}
\kn{ಹಿಂಗಾರು} \kn{ಮಳೆ} \kn{ಮುಂಗಾರಿನ} \kn{ನಂತರ} \kn{ಬರುತ್ತದೆ}.
\end{quote}

\begin{quote}
\emph{hiṅgāru maḷe muṅgārina nantara baruttade.}
\end{quote}

\begin{quote}
The late monsoon rain comes after the early monsoon.
\end{quote}

Its partner is \kn{ಮುಂಗಾರು}, the early rains.

\begin{cousinweb}
\textbf{\kn{ಹಿಂಗಾರು}} (\emph{hiṅgāru}) is a native Kannada compound of \emph{pin}, \emph{him}, \textquotedblleft{}behind, later\textquotedblright{} (DEDR 4205), and \emph{kār}, \textquotedblleft{}the rainy season\textquotedblright{} (DEDR 1278).
\end{cousinweb}

\begin{grammarlens}[title={What you've built}]
One more word for getting around.
\end{grammarlens}

\subsection*{Your turn}
\begin{itemize}
\raggedright
  \item \emph{Say it:} \emph{hiṅgāru}
  \item \emph{Say it:} \emph{hiṅgāru}, once more
  \item \emph{Say it:} \emph{hiṅgāru maḷe muṅgārina nantara baruttade.} — The late monsoon rain comes after the early monsoon.
  \item \emph{Recall:} say the Kannada for a downpour, then the Kannada for to get soaked, then say \emph{hiṅgāru} again
\end{itemize}

\begin{tcolorbox}[breakable,colback=teal!4,colframe=teal!35!black,title={Before you move on}]
What does \kn{ಹಿಂಗಾರು} mean? (\textquotedblleft{}The late monsoon; the retreating rains\textquotedblright{}.) Say it once more.
\end{tcolorbox}
\section[\texorpdfstring{\kn{ಆಲಿಕಲ್ಲು} (ālikallu)}{ಆಲಿಕಲ್ಲು (ālikallu)}]{\kn{ಆಲಿಕಲ್ಲು} (ālikallu) — hail; a hailstone}

\label{lesson:KA-C486-alikallu}



\begin{quote}
Before the new one: say the Kannada for to get soaked, then the Kannada for the late monsoon.
\end{quote}

\subsection*{You'll want to know: \kn{ಆಲಿಕಲ್ಲು}}
\textbf{\kn{ಆಲಿಕಲ್ಲು}} — \emph{ālikallu} — \textquotedblleft{}hail; a hailstone\textquotedblright{}.

\begin{quote}
\kn{ಬಿರುಗಾಳಿಯ} \kn{ಜೊತೆಗೆ} \kn{ಆಲಿಕಲ್ಲು} \kn{ಬಿದ್ದಿತು}.
\end{quote}

\begin{quote}
\emph{birugāḷiya jotege ālikallu bidditu.}
\end{quote}

\begin{quote}
Hail fell along with the storm.
\end{quote}

\kn{ಆಲಿ}, hail, plus \kn{ಕಲ್ಲು}, a stone: stones of ice.

\begin{cousinweb}
\textbf{\kn{ಆಲಿಕಲ್ಲು}} (\emph{ālikallu}) is native Dravidian: DEDR 384 lists Kannada \emph{āli-kal}, \emph{āle-kallu}, \textquotedblleft{}a hailstone\textquotedblright{}, built with \kn{ಕಲ್ಲು}, \textquotedblleft{}a stone\textquotedblright{} (DEDR 1298).
\end{cousinweb}

\begin{grammarlens}[title={What you've built}]
One more word for getting around.
\end{grammarlens}

\subsection*{Your turn}
\begin{itemize}
\raggedright
  \item \emph{Say it:} \emph{ālikallu}
  \item \emph{Say it:} \emph{ālikallu}, once more
  \item \emph{Say it:} \emph{birugāḷiya jotege ālikallu bidditu.} — Hail fell along with the storm.
  \item \emph{Recall:} say the Kannada for to get soaked, then the Kannada for the late monsoon, then say \emph{ālikallu} again
\end{itemize}

\begin{tcolorbox}[breakable,colback=teal!4,colframe=teal!35!black,title={Before you move on}]
What does \kn{ಆಲಿಕಲ್ಲು} mean? (\textquotedblleft{}Hail; a hailstone\textquotedblright{}.) Say it once more.
\end{tcolorbox}
\section[\texorpdfstring{\kn{ಮಬ್ಬು} (mabbu)}{ಮಬ್ಬು (mabbu)}]{\kn{ಮಬ್ಬು} (mabbu) — haze; dimness}

\label{lesson:KA-C486-mabbu}



\begin{quote}
Before the new one: say the Kannada for the late monsoon, then the Kannada for hail.
\end{quote}

\subsection*{You'll want to know: \kn{ಮಬ್ಬು}}
\textbf{\kn{ಮಬ್ಬು}} — \emph{mabbu} — \textquotedblleft{}haze; dimness\textquotedblright{}.

\begin{quote}
\kn{ಬೆಳಗ್ಗೆ} \kn{ನಗರದ} \kn{ಮೇಲೆ} \kn{ಮಬ್ಬು} \kn{ಇತ್ತು}.
\end{quote}

\begin{quote}
\emph{beḷagge nagarada mēle mabbu ittu.}
\end{quote}

\begin{quote}
In the morning there was a haze over the city.
\end{quote}

Also the dimness of dusk: \kn{ಮಬ್ಬು} \kn{ಬೆಳಕು}, dim light.

\begin{cousinweb}
\textbf{\kn{ಮಬ್ಬು}} (\emph{mabbu}) is native Dravidian: DEDR 4728 lists Kannada \emph{marvu}, \emph{mabbu}, \textquotedblleft{}darkness, gloom, dimness\textquotedblright{}.
\end{cousinweb}

\begin{grammarlens}[title={What you've built}]
One more word for getting around.
\end{grammarlens}

\subsection*{Your turn}
\begin{itemize}
\raggedright
  \item \emph{Say it:} \emph{mabbu}
  \item \emph{Say it:} \emph{mabbu}, once more
  \item \emph{Say it:} \emph{beḷagge nagarada mēle mabbu ittu.} — In the morning there was a haze over the city.
  \item \emph{Recall:} say the Kannada for the late monsoon, then the Kannada for hail, then say \emph{mabbu} again
\end{itemize}

\begin{tcolorbox}[breakable,colback=teal!4,colframe=teal!35!black,title={Before you move on}]
What does \kn{ಮಬ್ಬು} mean? (\textquotedblleft{}Haze; dimness\textquotedblright{}.) Say it once more.
\end{tcolorbox}
\section[\texorpdfstring{\kn{ಮುಸುಕು} (musuku)}{ಮುಸುಕು (musuku)}]{\kn{ಮುಸುಕು} (musuku) — to shroud; to cover over}

\label{lesson:KA-C486-musuku}



\begin{quote}
Before the new one: say the Kannada for hail, then the Kannada for haze.
\end{quote}

\subsection*{You'll want to know: \kn{ಮುಸುಕು}}
\textbf{\kn{ಮುಸುಕು}} — \emph{musuku} — \textquotedblleft{}to shroud; to cover over\textquotedblright{}.

\begin{quote}
\kn{ಬೆಳಗ್ಗೆ} \kn{ಬೆಟ್ಟಗಳನ್ನು} \kn{ಮಂಜು} \kn{ಮುಸುಕಿತು}.
\end{quote}

\begin{quote}
\emph{beḷagge beṭṭagaḷannu mañju musukitu.}
\end{quote}

\begin{quote}
In the morning mist shrouded the hills.
\end{quote}

Past (it …): \kn{ಮುಸುಕಿತು}, \emph{musukitu}.

As a noun, a veil: \kn{ಮುಸುಕು} \kn{ಹಾಕು}, to cover one's head.

\begin{cousinweb}
\textbf{\kn{ಮುಸುಕು}} (\emph{musuku}) is native Dravidian: DEDR 5030 lists Kannada \emph{musuku}, \emph{musugu}, \textquotedblleft{}to cover, to hide\textquotedblright{}.
\end{cousinweb}

\begin{grammarlens}[title={What you've built}]
One more word for getting around.
\end{grammarlens}

\subsection*{Your turn}
\begin{itemize}
\raggedright
  \item \emph{Say it:} \emph{musuku}
  \item \emph{Say it:} \emph{musuku}, once more
  \item \emph{Say it:} \emph{beḷagge beṭṭagaḷannu mañju musukitu.} — In the morning mist shrouded the hills.
  \item \emph{Recall:} say the Kannada for hail, then the Kannada for haze, then say \emph{musuku} again
\end{itemize}

\begin{tcolorbox}[breakable,colback=teal!4,colframe=teal!35!black,title={Before you move on}]
What does \kn{ಮುಸುಕು} mean? (\textquotedblleft{}To shroud; to cover over\textquotedblright{}.) Say it once more.
\end{tcolorbox}
\section[\texorpdfstring{\kn{ಹಿಮಪಾತ} (himapāta)}{ಹಿಮಪಾತ (himapāta)}]{\kn{ಹಿಮಪಾತ} (himapāta) — snowfall}

\label{lesson:KA-C486-himapata}



\begin{quote}
Before the new one: say the Kannada for haze, then the Kannada for to shroud.
\end{quote}

\subsection*{You'll want to know: \kn{ಹಿಮಪಾತ}}
\textbf{\kn{ಹಿಮಪಾತ}} — \emph{himapāta} — \textquotedblleft{}snowfall\textquotedblright{}.

\begin{quote}
\kn{ಹಿಮಪಾತದಿಂದ} \kn{ಬೆಟ್ಟದ} \kn{ರಸ್ತೆ} \kn{ಮುಚ್ಚಿತು}.
\end{quote}

\begin{quote}
\emph{himapātadinda beṭṭada raste muccitu.}
\end{quote}

\begin{quote}
The mountain road was closed by snowfall.
\end{quote}

\kn{ಹಿಮ}, snow, plus \kn{ಪಾತ}, a fall.

\begin{cousinweb}
\textbf{\kn{ಹಿಮಪಾತ}} (\emph{himapāta}) is built from two Sanskrit words: \emph{hima}, \textquotedblleft{}cold, frost, snow\textquotedblright{} (CDIAL 14096), and \emph{pāta}, \textquotedblleft{}a fall\textquotedblright{} (CDIAL 8052).
\end{cousinweb}

\begin{grammarlens}[title={What you've built}]
One more word for getting around.
\end{grammarlens}

\subsection*{Your turn}
\begin{itemize}
\raggedright
  \item \emph{Say it:} \emph{himapāta}
  \item \emph{Say it:} \emph{himapāta}, once more
  \item \emph{Say it:} \emph{himapātadinda beṭṭada raste muccitu.} — The mountain road was closed by snowfall.
  \item \emph{Recall:} say the Kannada for haze, then the Kannada for to shroud, then say \emph{himapāta} again
\end{itemize}

\begin{tcolorbox}[breakable,colback=teal!4,colframe=teal!35!black,title={Before you move on}]
What does \kn{ಹಿಮಪಾತ} mean? (\textquotedblleft{}Snowfall\textquotedblright{}.) Say it once more.
\end{tcolorbox}
